\documentclass[a4paper,12pt]{ctexart}
\usepackage[left=2.5cm, right=2.5cm, top=2.5cm, bottom=2.5cm]{geometry}
\usepackage{amsmath}
\usepackage{amssymb}
\usepackage{booktabs} % 优美的三线表
\usepackage{array}
\usepackage{graphicx}
\usepackage{float}
\usepackage{caption}  % 更好的标题控制
\usepackage{enumitem} % 列表定制
\usepackage{amsthm}   % 定理环境

% --- 视觉与排版优化设置 ---

% 1. 增加行距，中文文档通常需要 1.3-1.5 倍行距才舒适
\linespread{1.4}

% 2. 设置段落间距
\setlength{\parskip}{0.6em}
\setlength{\parindent}{2em}

% 3. 增加多行公式之间的间距
\setlength{\jot}{8pt}

% 4. 表格行高设置（防止表格内容挤在一起）
\renewcommand{\arraystretch}{1.3}

% 5. 自定义环境：让“例”和“解”看起来更专业
\newtheoremstyle{exampleStyle}
  {1em}   % 上方间距
  {1em}   % 下方间距
  {\normalfont} % 主体字体
  {0pt}   % 缩进
  {\bfseries}   % 头部字体加粗
  {}      % 头部后标点
  {1em}   % 头部后间距
  {}      % 头部格式
\theoremstyle{exampleStyle}
\newtheorem*{example}{例}

% 定义“解”的样式，不用环境，直接用命令更灵活
\newcommand{\solution}{\par\noindent\textbf{解}\quad}
\newcommand{\prooftext}{\par\noindent\textbf{证明}\quad}

% --- 文档开始 ---

\begin{document}

% Page 1
\section*{第一章}

\begin{example}[3]
要使 $\sqrt{20}$ 的近似值的相对误差限小于 $0.1\%$，要取几位有效数字？
\end{example}

\solution 设取 $n$ 位有效数字，由定理 1，$\varepsilon_r^* \leqslant \frac{1}{2a_1} \times 10^{-n+1}$。由于 $\sqrt{20} = 4.4 \cdots$，知 $a_1=4$，故只要取 $n=4$，就有
\[
\varepsilon_r^* \leqslant 0.125 \times 10^{-3} < 10^{-3} = 0.1\%,
\]
即只要对 $\sqrt{20}$ 的近似值取 4 位有效数字，其相对误差限就小于 $0.1\%$。此时由开方表得 $\sqrt{20} \approx 4.472$。

\begin{example}[4]
已测得某场地长 $l$ 的值为 $l^* = 110\,\mathrm{m}$，宽 $d$ 的值为 $d^* = 80\,\mathrm{m}$，已知 $|l - l^*| \leqslant 0.2\,\mathrm{m}, |d - d^*| \leqslant 0.1\,\mathrm{m}$。试求面积 $s = ld$ 的绝对误差限与相对误差限。
\end{example}

\solution 因 $s=ld$，$\frac{\partial s}{\partial l} = d$，$\frac{\partial s}{\partial d} = l$，由 (2.3) 式知
\[
\varepsilon(s^*) \approx \left| \left( \frac{\partial s}{\partial l} \right)^* \right| \varepsilon(l^*) + \left| \left( \frac{\partial s}{\partial d} \right)^* \right| \varepsilon(d^*),
\]
其中
\[
\left( \frac{\partial s}{\partial l} \right)^* = d^* = 80\,\mathrm{m}, \quad \left( \frac{\partial s}{\partial d} \right)^* = l^* = 110\,\mathrm{m},
\]
而 $\varepsilon(l^*) = 0.2\,\mathrm{m}, \varepsilon(d^*) = 0.1\,\mathrm{m}$，于是绝对误差限
\[
\varepsilon(s^*) \approx 80 \times 0.2 + 110 \times 0.1 = 27\,(\mathrm{m}^2);
\]
相对误差限
\[
\varepsilon_r(s^*) = \frac{\varepsilon(s^*)}{|s^*|} = \frac{\varepsilon(s^*)}{l^* d^*} \approx \frac{27}{8800} = 0.31\%.
\]

\section*{第二章}

% Page 2
\begin{example}[2]
已给 $\sin 0.32 = 0.314567, \sin 0.34 = 0.333487, \sin 0.36 = 0.352274$，用线性插值及抛物插值计算 $\sin 0.3367$ 的值并估计截断误差。
\end{example}

\solution 由题意取 $x_0 = 0.32, y_0 = 0.314567, x_1 = 0.34, y_1 = 0.333487, x_2 = 0.36, y_2 = 0.352274$。

用线性插值计算，由于 $0.3367$ 介于 $x_0, x_1$ 之间，故取 $x_0, x_1$ 进行计算，由公式 (2.1) 得
\begin{align*}
\sin 0.3367 &\approx L_1(0.3367) = y_0 + \frac{y_1 - y_0}{x_1 - x_0}(0.3367 - x_0) \\
&= 0.314567 + \frac{0.01892}{0.02} \times 0.0167 = 0.330365.
\end{align*}
由 (2.15) 式得其截断误差
\[
|R_1(x)| \leqslant \frac{M_2}{2} \left| (x - x_0)(x - x_1) \right|,
\]
其中 $M_2 = \max_{x_0 \leqslant x \leqslant x_1} |f''(x)|$。因 $f(x) = \sin x, f''(x) = -\sin x$，可取 $M_2 = \max_{x_0 \leqslant x \leqslant x_1} |\sin x| = \sin x_1 \leqslant 0.3335$，于是
\begin{align*}
|R_1(0.3367)| &= |\sin 0.3367 - L_1(0.3367)| \\
&\leqslant \frac{1}{2} \times 0.3335 \times 0.0167 \times 0.0033 \leqslant 0.92 \times 10^{-5}.
\end{align*}

用抛物线插值计算 $\sin 0.3367$ 时，由公式 (2.5) 得
\begin{align*}
\sin 0.3367 &\approx y_0 \frac{(x - x_1)(x - x_2)}{(x_0 - x_1)(x_0 - x_2)} + y_1 \frac{(x - x_0)(x - x_2)}{(x_1 - x_0)(x_1 - x_2)} + y_2 \frac{(x - x_0)(x - x_1)}{(x_2 - x_0)(x_2 - x_1)} \\
&= L_2(0.3367) \\
&= 0.314567 \times \frac{0.7689 \times 10^{-4}}{0.0008} + 0.333487 \times \frac{3.89 \times 10^{-4}}{0.0004} \\
&\quad + 0.352274 \times \frac{-0.5511 \times 10^{-4}}{0.0008} \\
&= 0.330374.
\end{align*}
这个结果与 6 位有效数字的正弦函数表完全一样，这说明查表时用二次插值精度已相当高了。由 (2.14) 式得其截断误差限
\[
|R_2(x)| \leqslant \frac{M_3}{6} \left| (x - x_0)(x - x_1)(x - x_2) \right|,
\]
其中
\[
M_3 = \max_{x_0 \leqslant x \leqslant x_2} |f'''(x)| = \cos x_0 < 0.9493,
\]
于是
\begin{align*}
|R_2(0.3367)| &= |\sin 0.3367 - L_2(0.3367)| \\
&\leqslant \frac{1}{6} \times 0.9493 \times 0.0167 \times 0.033 \times 0.0233 < 2.0132 \times 10^{-6}.
\end{align*}

% Page 3
\begin{example}[4]
给出 $f(x)$ 的函数表（见表 2-2），求 4 次牛顿插值多项式，并由此计算 $f(0.596)$ 的近似值。
\end{example}

首先根据给定函数表造出均差表。

\begin{table}[H]
\centering
\caption*{表 2-2 函数及均差表}
% 使用 p{...} 或适当间距来防止表格过宽
\begin{tabular}{c c c c c c c}
\toprule
$x$ & $f(x)$ & 一阶 & 二阶 & 三阶 & 四阶 & 五阶 \\
\midrule
$0.40$ & $0.41075$ & & & & & \\
$0.55$ & $0.57815$ & $1.11600$ & & & & \\
$0.65$ & $0.69675$ & $1.18600$ & $0.28000$ & & & \\
$0.80$ & $0.88811$ & $1.27573$ & $0.35893$ & $0.19733$ & & \\
$0.90$ & $1.02652$ & $1.38410$ & $0.43348$ & $0.21300$ & $0.03134$ & \\
$1.05$ & $1.25382$ & $1.51533$ & $0.52493$ & $0.22863$ & $0.03126$ & $-0.00012$ \\
\bottomrule
\end{tabular}
\end{table}

从均差表看到 4 阶均差近似常数，故取 4 次插值多项式 $P_4(x)$ 做近似即可。
\begin{align*}
P_4(x) &= 0.41075 + 1.116(x - 0.4) + 0.28(x - 0.4)(x - 0.55) \\
&\quad + 0.19733(x - 0.4)(x - 0.55)(x - 0.65) \\
&\quad + 0.03134(x - 0.4)(x - 0.55)(x - 0.65)(x - 0.8),
\end{align*}
于是
\[
f(0.596) \approx P_4(0.596) = 0.63192,
\]
截断误差
\[
|R_4(x)| \approx |f[x_0, x_1, \cdots, x_5] \omega_5(0.596)| \leqslant 3.63 \times 10^{-9}.
\]
这说明截断误差很小，可忽略不计。

此例的截断误差估计中，5 阶均差 $f[x_0, \cdots, x_4]$ 用 $f[x_0, x_1, \cdots, x_5] = -0.00012$ 近似。另一种方法是取 $x = 0.596$，由 $f(0.596) \approx 0.63192$，可求得 $f[x, x_0, \cdots, x_4]$ 的近似值，从而可求得 $|R_4(x)|$ 的近似值。

\section*{第三章}

% Page 4
\begin{example}[9]
已知一组实验数据如表 3-1，求它的拟合曲线。
\end{example}

\begin{table}[H]
\centering
\caption*{表 3-1 实验数据}
\begin{tabular}{c c c c c c}
\toprule
$x_i$ & $1$ & $2$ & $3$ & $4$ & $5$ \\
\midrule
$f_i$ & $4$ & $4.5$ & $6$ & $8$ & $8.5$ \\
$\omega_i$ & $2$ & $1$ & $3$ & $1$ & $1$ \\
\bottomrule
\end{tabular}
\end{table}

\solution 根据所给数据，在坐标纸上标出，见图 3-5。从图中看到各点在一条直线附近，故可选择线性函数作拟合曲线，即令 $S_1^*(x) = a_0 + a_1 x$，这里 $m=4, n=1, \varphi_0(x) = 1, \varphi_1(x) = x$，故
\begin{gather*}
(\varphi_0, \varphi_0) = \sum_{i=0}^4 \omega_i = 8, \\
(\varphi_0, \varphi_1) = (\varphi_1, \varphi_0) = \sum_{i=0}^4 \omega_i x_i = 22, \\
(\varphi_1, \varphi_1) = \sum_{i=0}^4 \omega_i x_i^2 = 74, \quad (\varphi_0, f) = \sum_{i=0}^4 \omega_i f_i = 47, \\
(\varphi_1, f) = \sum_{i=0}^4 \omega_i x_i f_i = 145.5.
\end{gather*}
由法方程 (4.6) 得线性方程组
\[
\begin{cases}
8a_0 + 22a_1 = 47, \\
22a_0 + 74a_1 = 145.5.
\end{cases}
\]
解得 $a_0 = 2.5648, a_1 = 1.2037$。于是所求拟合曲线为
\[
S_1^*(x) = 2.5648 + 1.2037x.
\]

\begin{example}[10]
设数据 $(x_i, y_i) \, (i=0, 1, 2, 3, 4)$ 由表 3-2 给出，表中第 4 行为 $\ln y_i = \bar{y}_i$，可以看出数学模型为 $y = ae^{bx}$，用最小二乘法确定 $a$ 及 $b$。
\end{example}

\solution 根据给定数据 $(x_i, y_i)$ 描图可确定拟合曲线方程为 $y = ae^{bx}$，它不是线性形式。两边取对数得 $\ln y = \ln a + bx$，若令 $\bar{y} = \ln y, A = \ln a$，则得 $\bar{y} = A + bx, \varphi = \{1, x\}$。为确定 $A, b$，先将 $(x_i, y_i)$ 转化为 $(x_i, \bar{y}_i)$。

根据最小二乘法，取 $\varphi_0(x) = 1, \varphi_1(x) = x, \omega(x) \equiv 1$，得

\begin{table}[H]
\centering
\caption*{表 3-2 数据表}
\begin{tabular}{c c c c c c}
\toprule
$i$ & $0$ & $1$ & $2$ & $3$ & $4$ \\
\midrule
$x_i$ & $1.00$ & $1.25$ & $1.50$ & $1.75$ & $2.00$ \\
$y_i$ & $5.10$ & $5.79$ & $6.53$ & $7.45$ & $8.46$ \\
$\bar{y}_i$ & $1.629$ & $1.756$ & $1.876$ & $2.008$ & $2.135$ \\
\bottomrule
\end{tabular}
\end{table}

\begin{gather*}
(\varphi_0, \varphi_0) = 5, \quad (\varphi_0, \varphi_1) = \sum_{i=0}^4 x_i = 7.5, \quad (\varphi_1, \varphi_1) = \sum_{i=0}^4 x_i^2 = 11.875, \\
(\varphi_0, \bar{y}) = \sum_{i=0}^4 \bar{y}_i = 9.404, \quad (\varphi_1, \bar{y}) = \sum_{i=0}^4 x_i \bar{y}_i = 14.422.
\end{gather*}
故有法方程
\[
\begin{cases}
5A + 7.50b = 9.404, \\
7.50A + 11.875b = 14.422.
\end{cases}
\]
解得 $A = 1.122, b = 0.505, a = e^A = 3.071$。于是得最小二乘拟合曲线为
\[
y = 3.071e^{0.505x}.
\]
现在很多数学软件配有自动选择数学模型的程序，其方法与本例同。

\section*{第四章}

\begin{example}[1]
给定形如 $\int_0^1 f(x)dx \approx A_0 f(0) + A_1 f(1) + B_0 f'(0)$ 的求积公式，试确定系数 $A_0, A_1, B_0$，使公式具有尽可能高的代数精度。
\end{example}

\solution 根据题意可令 $f(x) = 1, x, x^2$ 分别代入求积公式使它精确成立：
\begin{align*}
\text{当 } f(x)=1 \text{ 时，} & A_0 + A_1 = \int_0^1 1 \cdot dx = 1; \\
\text{当 } f(x)=x \text{ 时，} & A_1 + B_0 = \int_0^1 x dx = \frac{1}{2}; \\
\text{当 } f(x)=x^2 \text{ 时，} & A_1 = \int_0^1 x^2 dx = \frac{1}{3}.
\end{align*}
解得 $A_1 = \frac{1}{3}, A_0 = \frac{2}{3}, B_0 = \frac{1}{6}$，于是有
\[
\int_0^1 f(x) dx \approx \frac{2}{3} f(0) + \frac{1}{3} f(1) + \frac{1}{6} f'(0).
\]
当 $f(x) = x^3$ 时 $\int_0^1 x^3 dx = \frac{1}{4}$。而上式右端为 $\frac{1}{3}$，故公式对 $f(x)=x^3$ 不精确成立，其代数精度为 2。

\begin{example}[3]
对于函数 $f(x) = \frac{\sin x}{x}$，给出 $n=8$ 时的函数表（见表 4-2），试用复合梯形公式 (3.2) 及复合辛普森公式 (3.5) 计算积分
$I = \int_0^1 \frac{\sin x}{x} dx$, 并估计误差。
\end{example}

\begin{table}[H]
\centering
\caption*{表 4-2 计算结果}
\begin{tabular}{c c | c c}
\toprule
$x$ & $f(x)$ & $x$ & $f(x)$ \\
\midrule
$0$ & $1$ & $5/8$ & $0.9361556$ \\
$1/8$ & $0.9973978$ & $3/4$ & $0.9088516$ \\
$1/4$ & $0.9896158$ & $7/8$ & $0.8771925$ \\
$3/8$ & $0.9767267$ & $1$ & $0.8414709$ \\
$1/2$ & $0.9588510$ & & \\
\bottomrule
\end{tabular}
\end{table}

\solution 将积分区间 $[0, 1]$ 划分为 8 等份，应用复合梯形法求得 $T_8 = 0.9456909$;
而如果将 $[0, 1]$ 分为 4 等份，应用复合辛普森法有 $S_4 = 0.9460832$.

比较上面两个结果 $T_8$ 与 $S_4$，它们都需要提供 9 个点上的函数值，计算量基本相同，然而精度却差别很大，同积分的准确值 $I = 0.9460831$ 比较，复合梯形法的结果 $T_8$ 只有两位有效数字，而复合辛普森法的结果 $S_4$ 却有六位有效数字。

为了利用余项公式估计误差，要求 $f(x) = \frac{\sin x}{x}$ 的高阶导数。由于
\[
f(x) = \frac{\sin x}{x} = \int_0^1 \cos(xt) dt,
\]
所以有
\[
f^{(k)}(x) = \int_0^1 \frac{d^k}{dx^k} (\cos xt) dt = \int_0^1 t^k \cos \left( xt + \frac{k\pi}{2} \right) dt,
\]
于是
\[
\max_{0 \leqslant x \leqslant 1} |f^{(k)}(x)| \leqslant \int_0^1 \left| \cos \left( xt + \frac{k\pi}{2} \right) \right| t^k dt \leqslant \int_0^1 t^k dt = \frac{1}{k+1}.
\]
由 (3.3) 式得复合梯形公式的误差
\[
|R_8(f)| = |I - T_n| \leqslant \frac{h^2}{12} \max_{0 \leqslant x \leqslant 1} |f''(x)| \leqslant \frac{1}{12} \left( \frac{1}{8} \right)^2 \frac{1}{3} = 0.000434.
\]
对复合辛普森公式的误差，由 (3.6) 式得
\[
|R_4(f)| = |I - S_4| \leqslant \frac{1}{2880} \left( \frac{1}{4} \right)^4 \frac{1}{5} = 0.271 \times 10^{-6}.
\]

\section*{第五章}

\begin{example}[1]
求 $A = \begin{pmatrix} 1 & -2 & 2 \\ -2 & -2 & 4 \\ 2 & 4 & -2 \end{pmatrix}$ 的特征值及谱半径。
\end{example}

\solution $A$ 的特征方程为
\begin{align*}
\det(\lambda I - A) &= \begin{vmatrix} \lambda - 1 & 2 & -2 \\ 2 & \lambda + 2 & -4 \\ -2 & -4 & \lambda + 2 \end{vmatrix} \\
&= \lambda^3 + 3\lambda^2 - 24\lambda + 28 \\
&= (\lambda - 2)^2 (\lambda + 7) = 0,
\end{align*}
故 $A$ 的特征值为 $\lambda_1 = \lambda_2 = 2, \lambda_3 = 7, A$ 的谱半径为 $\rho(A) = 7$。

\begin{example}[5]
用直接三角分解法解
$\begin{pmatrix} 1 & 2 & 3 \\ 2 & 5 & 2 \\ 3 & 1 & 5 \end{pmatrix} \begin{pmatrix} x_1 \\ x_2 \\ x_3 \end{pmatrix} = \begin{pmatrix} 14 \\ 18 \\ 20 \end{pmatrix}$.
\end{example}

\solution 用分解公式 (3.2), (3.3) 计算得
\[
A = \begin{pmatrix} 1 & 2 & 3 \\ 2 & 5 & 2 \\ 3 & 1 & 5 \end{pmatrix} \xrightarrow{} \begin{pmatrix} 1 & 0 & 0 \\ 2 & 1 & 0 \\ 3 & -5 & 1 \end{pmatrix} \begin{pmatrix} 1 & 2 & 3 \\ 0 & 1 & -4 \\ 0 & 0 & -24 \end{pmatrix} = LU.
\]
求解
\begin{align*}
Ly &= (14, 18, 20)^T, \quad \text{得 } y = (14, -10, -72)^T, \\
Ux &= (14, -10, -72)^T, \quad \text{得 } x = (1, 2, 3)^T.
\end{align*}

\vspace{1em}
\noindent\textbf{2. 选主元的三角分解法}

\textbf{算法 2}（选主元的三角分解法） 设 $Ax = b$，其中 $A$ 为非奇异矩阵。本算法采用选主元的三角分解法，用 $PA = I_{n-1} \cdots I_1 A$ 的三角分解冲掉 $A$，用整型数组 Ip($n$) 记录主行，解 $x$ 存放在 $b$ 内。

\begin{enumerate}[label=\textbf{\arabic*.}]
    \item 对于 $r = 1, 2, \cdots, n$
    \begin{enumerate}[label=(\arabic*)]
        \item 计算 $s_i$：
        $a_{ir} \leftarrow s_i = a_{ir} - \sum_{k=1}^{r-1} l_{ik} u_{kr}, \quad i = r, r+1, \cdots, n$
        \item 选主元：
        $|s_{i_r}| = \max_{r \leqslant i \leqslant n} |s_i|, \quad \text{Ip}(r) \leftarrow i$
        \item 交换 $A$ 的 $r$ 行与 $i_r$ 行元素：
        $a_{ri} \longleftrightarrow a_{i_r i}, \quad i = 1, 2, \cdots, n$
        \item 计算 $U$ 的第 $r$ 行元素，$L$ 的第 $r$ 列元素：
        $a_{rr} = u_{rr} = s_r$
        \[
        a_{ir} \leftarrow l_{ir} = s_i / u_{rr} = a_{ir} / a_{rr}, \quad i = r+1, \cdots, n, \text{且 } r \neq n
        \]
        \[
        a_{ri} \leftarrow u_{ri} = a_{ri} - \sum_{k=1}^{r-1} l_{rk} u_{ki}, \quad i = r+1, \cdots, n, \text{且 } r \neq n
        \]
        （这时有 $|l_{ir}| \leqslant 1$）
    \end{enumerate}

    \item 求解 $Ly = Pb$ 及 $Ux = y$：
    对于 $i = 1, 2, \cdots, n-1$
    \begin{enumerate}[label=(\arabic*)]
        \item $t \leftarrow \text{Ip}(i)$
        \item 如果 $i=t$ 则转(3)，否则 $b_i \longleftrightarrow b_t$
        \item （继续循环）
    \end{enumerate}
    \item $b_i \leftarrow b_i - \sum_{k=1}^{i-1} l_{ik} b_k, i = 2, 3, \cdots, n$
    \item $b_n \leftarrow b_n / u_{nn}, b_i \leftarrow (b_i - \sum_{k=i+1}^n u_{ik} b_k) / u_{ii}, i = n-1, \cdots, 1$
\end{enumerate}

\section*{第六章}

\begin{example}[8]
在线性方程组 $Ax = b$ 中，
$A = \begin{pmatrix} 1 & a & a \\ a & 1 & a \\ a & a & 1 \end{pmatrix}$,
证明当 $-\frac{1}{2} < a < 1$ 时高斯-塞德尔法收敛，而雅可比迭代法只在 $-\frac{1}{2} < a < \frac{1}{2}$ 时才收敛。
\end{example}

\prooftext 只要证 $-\frac{1}{2} < a < 1$ 时 $A$ 正定，由 $A$ 的顺序主子式 $\Delta_2 = \begin{vmatrix} 1 & a \\ a & 1 \end{vmatrix} = 1 - a^2 > 0$，得 $|a| < 1$，而 $\Delta_3 = \det A = 1 + 2a^3 - 3a^2 = (1 - a)^2 (1 + 2a) > 0$，得 $a > -\frac{1}{2}$，于是得到 $-\frac{1}{2} < a < 1$ 时 $\Delta_1 > 0, \Delta_2 > 0, \Delta_3 > 0$，故 $A$ 正定，高斯-塞德尔法收敛。

对雅可比迭代矩阵
\[
J = \begin{pmatrix} 0 & -a & -a \\ -a & 0 & -a \\ -a & -a & 0 \end{pmatrix},
\]
有 $\det(\lambda I - J) = \lambda^3 - 3\lambda a^2 + 2a^3 = (\lambda - a)^2 (\lambda + 2a) = 0$。
当 $\rho(J) = |2a| < 1$，即 $|a| < \frac{1}{2}$ 时雅可比法收敛。

\section*{第七章}

\begin{example}[2]
求方程 $f(x) = x^3 - x - 1 = 0$ 在区间 $[1.0, 1.5]$ 内的一个实根，要求准确到小数点后的第 2 位。
\end{example}

\solution 这里 $a=1.0, b=1.5$，而 $f(a) < 0, f(b) > 0$。取 $[a, b]$ 的中点 $x_0 = 1.25$，将区间二等分，由于 $f(x_0) < 0$，即 $f(x_0)$ 与 $f(a)$ 同号，故所求的根 $x^*$ 必在 $x_0$ 右侧，这时应令 $a_1 = x_0 = 1.25, b_1 = b = 1.5$，从而得到新的有根区间 $[a_1, b_1]$。

\begin{table}[H]
\centering
\caption*{表 7-2 计算结果}
\begin{tabular}{c c c c c}
\toprule
$k$ & $a_k$ & $b_k$ & $x_k$ & $f(x_k)$符号 \\
\midrule
0 & 1.0 & 1.5 & 1.25 & $-$ \\
1 & 1.25 & & 1.375 & $+$ \\
2 & & 1.375 & 1.3125 & $-$ \\
3 & 1.3125 & & 1.3438 & $+$ \\
4 & & 1.3438 & 1.3281 & $+$ \\
5 & & 1.3281 & 1.3203 & $-$ \\
6 & 1.3203 & & 1.3242 & \\
\bottomrule
\end{tabular}
\end{table}

二分法计算步骤：
\begin{description}[style=multiline, leftmargin=3em]
    \item[步骤 1] \textbf{准备}：计算 $f(x)$ 在有根区间 $[a, b]$ 端点处的值 $f(a), f(b)$。
    \item[步骤 2] \textbf{二分}：计算 $f(x)$ 在区间中点 $\frac{a+b}{2}$ 处的值 $f\left(\frac{a+b}{2}\right)$。
    \item[步骤 3] \textbf{判断}：若 $f\left(\frac{a+b}{2}\right)=0$，则 $\frac{a+b}{2}$ 即是根；否则检验：若 $f\left(\frac{a+b}{2}\right)f(a) < 0$，则以 $\frac{a+b}{2}$ 代替 $b$，否则以 $\frac{a+b}{2}$ 代替 $a$。
\end{description}

\begin{example}[4]
用不同方法求方程 $x^2 - 3 = 0$ 的根 $x^* = \sqrt{3}$。
\end{example}

\solution 构造不同的迭代法：
\begin{enumerate}[label=(\arabic*)]
    \item $x_{k+1} = x_k^2 + x_k - 3, \quad \varphi'(x^*) = 2\sqrt{3} + 1 > 1$.
    \item $x_{k+1} = \frac{3}{x_k}, \quad \varphi'(x^*) = -1$.
    \item $x_{k+1} = x_k - \frac{1}{4}(x_k^2 - 3), \quad \varphi'(x^*) \approx 0.134 < 1$.
    \item $x_{k+1} = \frac{1}{2} \left( x_k + \frac{3}{x_k} \right), \quad \varphi'(x^*) = 0$.
\end{enumerate}

\begin{table}[H]
\centering
\caption*{表 7-4 计算结果}
\begin{tabular}{c c c c c c}
\toprule
$k$ & $x_k$ & 迭代法(1) & 迭代法(2) & 迭代法(3) & 迭代法(4) \\
\midrule
0 & $x_0$ & 2 & 2 & 2 & 2 \\
1 & $x_1$ & 3 & 1.5 & 1.75 & 1.75 \\
2 & $x_2$ & 9 & 2 & 1.73475 & 1.732143 \\
3 & $x_3$ & 87 & 1.5 & 1.732361 & 1.732051 \\
\bottomrule
\end{tabular}
\end{table}

\begin{example}[7]
用牛顿法解方程 $xe^x - 1 = 0$.
\end{example}

\solution 牛顿公式为
$x_{k+1} = x_k - \frac{x_k - e^{-x_k}}{1 + x_k}$. 取迭代初值 $x_0 = 0.5$，迭代结果如下：

\begin{table}[H]
\centering
\caption*{表 7-7 计算结果}
\begin{tabular}{c c}
\toprule
$k$ & $x_k$ \\
\midrule
0 & 0.5 \\
1 & 0.57102 \\
2 & 0.56716 \\
3 & 0.56714 \\
\bottomrule
\end{tabular}
\end{table}

\section*{第八章}

\begin{example}[1]
求解初值问题 $\begin{cases} y' = y - \frac{2x}{y}, \quad 0 < x < 1, \\ y(0) = 1. \end{cases}$
\end{example}

\solution 欧拉公式的具体形式为
$y_{n+1} = y_n + h \left( y_n - \frac{2x_n}{y_n} \right)$.
取步长 $h=0.1$，计算结果见表 9-1。

\begin{table}[H]
\centering
\caption*{表 9-1 计算结果对比}
\begin{tabular}{c c c | c c c}
\toprule
$x_n$ & $y_n$ & $y(x_n)$ & $x_n$ & $y_n$ & $y(x_n)$ \\
\midrule
0.1 & 1.1000 & 1.0954 & 0.6 & 1.5090 & 1.4832 \\
0.2 & 1.1918 & 1.1832 & 0.7 & 1.5803 & 1.5492 \\
0.3 & 1.2774 & 1.2649 & 0.8 & 1.6498 & 1.6125 \\
0.4 & 1.3582 & 1.3416 & 0.9 & 1.7178 & 1.6733 \\
0.5 & 1.4351 & 1.4142 & 1.0 & 1.7848 & 1.7321 \\
\bottomrule
\end{tabular}
\end{table}

为了分析计算公式的精度，通常用泰勒展开将 $y(x_{n+1})$ 在 $x_n$ 处展开：
\[
y(x_{n+1}) = y(x_n) + y'(x_n)h + \frac{h^2}{2} y''(\xi_n).
\]
可得欧拉法 (2.1) 的误差
\[
y(x_{n+1}) - y_{n+1} = \frac{h^2}{2} y''(\xi_n) \approx \frac{h^2}{2} y''(x_n).
\]

如果在 (2.4) 式中右端积分用右矩形公式 $hf(x_{n+1}, y(x_{n+1}))$ 近似，则得\textbf{后退的欧拉法}：
\[
y_{n+1} = y_n + hf(x_{n+1}, y_{n+1}). \eqno{(2.5)}
\]
隐式方程 (2.5) 通常用迭代法求解：
\[
y_{n+1}^{(k+1)} = y_n + hf(x_{n+1}, y_{n+1}^{(k)}), \quad k = 0, 1, \cdots. \eqno{(2.6)}
\]
由此可知，只要 $hL < 1$ 迭代法 (2.6) 就收敛到解 $y_{n+1}$。

\end{document}